% ============================================================
%  Section 2, Segment 2C — Alternate Tunings / 特殊调弦
%  原书页码 53-73
% ============================================================
\chapter{Alternate Tunings / 特殊调弦}
\label{ch:tunings}

\eng{Alternative tunings are when you tune your guitar to something other than standard tuning (EADGCE). They are a key feature of math rock guitar and are often favoured by bands that write texturally rich guitar music.}

\chn{特殊调弦是指将吉他调成与标准调弦（EADGCE）不同的音高。它们是数学摇滚吉他的一个关键特征，并且常被创作纹理丰富吉他音乐的乐队所喜爱。}

\paragraph{Examples of Alternate Tunings}
\eng{This section will focus on two of the most common alternative tunings used: FACGCE and DAEAC\#E (or CGDGBD, which is the same tuning two steps lower).}

\chn{本节将重点讨论两种最常用的特殊调弦：FACGCE 和 DAEAC\#E（或 CGDGBD，它其实是 DAEAC\#E 降两度后的同一调弦）。}

\eng{Some examples of FACGCE tuning in action include American Football – "Never Meant", Colossal – "Hot Probs", and TTNG – "Crocodile". Song examples in DAEAC\#E include Sleepy Dog – "Screamers", Algernon Cadwallader – "Some Kind of Cadwallader", and Tiny Moving Parts – "Caution".}

\chn{一些使用 FACGCE 调弦的实例包括 American Football 的《Never Meant》、Colossal 的《Hot Probs》以及 TTNG 的《Crocodile》。使用 DAEAC\#E 调弦的歌曲示例包括 Sleepy Dog 的《Screamers》、Algernon Cadwallader 的《Some Kind of Cadwallader》和 Tiny Moving Parts 的《Caution》。}

\paragraph{What Makes Certain Alternate Tunings Sound Great?}
\eng{You may be wondering “what makes some alternative tunings sound so damn nice?” The secret is in the note selection, and with a bit of simple music theory I can explain why.}

\chn{你可能会想“是什么让某些特殊调弦听起来如此悦耳？”秘密在于音符选择，通过一些简单的乐理我可以解释原因。}

\eng{If we look at the notes that make up standard tuning, E A D G B E, and the distance between the notes, we end up with a bunch of fourth intervals and one third interval. When strummed, this does not really sound like much. Technically, we could say the open strings of standard tuning form an Em11 chord.}

\chn{如果我们查看构成标准调弦的音符 E A D G B E 及其间距，我们会得到一连串的纯四度音程和一个大三度音程。扫弦时，它听起来并没有太特别。从技术上讲，我们可以说标准调弦的空弦形成了一个 Em11 和弦。}

\eng{If we take the tuning FACGCE and analyse it similarly, we will see why it sounds more like a chord than a random assortment of intervals. FACGCE contains more third intervals in its tuning, which makes it sound more like an Fmaj9 chord. This means the tuning is an open voiced chord tuning. It is based on a chord—in this case, an Fmaj9 chord. That is why it sounds so pleasing.}

\chn{如果我们用相同的方式分析 FACGCE 调弦，我们就会明白为什么它听起来更像一个和弦而不是随机间隔的音。FACGCE 的调弦中包含更多的大三度音程，使它听起来更像一个 Fmaj9 和弦。这意味着这种调弦是一个开放和弦调弦。它基于一个和弦——在此例中是 Fmaj9 和弦。这就是它听起来如此动听的原因。}

\eng{Now let's repeat the process for DAEAC\#E tuning. This tuning is a bit more challenging to work out, but it is still based around a chord: an A/D chord (an A major chord with a D root, also known as a slash chord). This tuning is slightly less complex harmonically than FACGCE, but that lower D note over an A chord sounds inspiring. The difference between these two tunings is that FACGCE contains a major seventh degree, which gives the tuning its characteristic happy-but-a-little-sad feel.}

\chn{现在我们对 DAEAC\#E 调弦重复这个过程。这个调弦稍微难以推导，但它仍然基于一个和弦：A/D 和弦（以 D 为根音的 A 大调和弦，也称作斜杠和弦）。与 FACGCE 相比，这个调弦在和声上略微简单一些，但低 D 音搭配 A 和弦听起来同样富有灵感。这两种调弦之间的区别在于 FACGCE 包含大七度音，这赋予了这种调弦一种快乐中带着一丝忧伤的特质。}

\paragraph{String Gauges for Alternate Tunings}
\eng{If you are only temporarily changing strings or adjusting certain strings by a semitone, it is usually fine to stick with the same string gauge you use for standard tuning. However, for most alternative tunings, you will need more significant changes. I suggest using your regular string gauge to experiment with the new tuning, and if you decide to use it more frequently, explore changing individual string gauges accordingly.}

\chn{如果你只是临时换弦或将某些弦调高/调低一个半音，通常仍然可以使用标准调弦时的弦规。然而，对于大多数特殊调弦，你将需要更明显的调整。我建议先用常用的弦规尝试新调弦，如果决定更频繁地使用它，再根据需要调整单根弦的弦规。}

\eng{If you plan to use a tuning for an extended period or make larger tuning adjustments, I recommend:}

\chn{如果你打算长时间使用某个调弦或进行较大的调弦变化，我建议：}

\begin{enumerate}[leftmargin=*]
  \item \eng{For strings tuned down, use a thicker gauge string, usually one gauge higher than your current string.}\chn{对于调低的弦，使用更粗的弦规，通常比当前的弦规大一号。}
  \item \eng{If you are tuning up, use a thinner gauge string.}\chn{如果你要调高弦，则使用更细的弦。}
  \item \eng{For tuning changes larger than a whole tone, experiment with thicker or thinner gauges depending on direction.}\chn{对于超过全音的调弦变化，根据调高或调低的方向尝试使用更粗或更细的弦。}
\end{enumerate}

\eng{This may require buying individual strings or hybrid sets, but it is worth it for better tension balance across the neck.}

\chn{这可能需要购买单支弦或混合弦组，但为了更均衡的琴颈张力，这是值得的。}

\paragraph{How to Write When Using Different Tunings?}
\eng{There is no easy answer or easy way of approaching writing in different tunings. I have a solid knowledge of where chords and notes are in standard tuning, but when I change to a different tuning, I lose that knowledge. You could look at this as off-putting, but I recommend seeing it as a good challenge.}

\chn{在不同调弦中创作没有简单的答案或捷径。我对标准调弦中和弦与音符的位置有扎实的认识，但当我换到不同的调弦时，我就会失去这种认识。你可能觉得这令人却步，但我建议把它当作一个好挑战。}

\eng{Exploring alternate tunings can also serve as excellent ear training. You will be challenged to find the notes you hear in your mind on the guitar, and this process can lead to discovering unique and appealing chords.}

\chn{探索特殊调弦也可以作为极好的听力训练。你会被挑战在吉他上找到你脑海中听到的音符，这个过程可能会引导你发现独特而动听的和弦。}

\eng{If you want to know where all the notes are in many different tunings, the tool https://fretboard-navigator.com/\# is very helpful. You can choose a tuning, a key, and even scales to see where they are in a tuning on the fretboard. Bookmark it if you have not already.}

\chn{如果你想知道不同调弦中所有音符的位置，工具 https://fretboard-navigator.com/\# 会非常有帮助。你可以选择调弦、调和音阶，查看它们在指板上的位置。如果你还没有收藏，建议收藏一下。}

\eng{Remember that open tunings offer a specific sound and feeling because they are tuned to a chord. You can use this to your advantage by incorporating the open tuning chord into your musical ideas and by including open strings within your chord choices. If you want to alter that feeling, use a capo or choose a different tuning instead.}

\chn{记住，开弦调弦之所以提供特定的声音和感觉，是因为它们调成了一个和弦。你可以利用这一点，把开弦调弦的和弦融入你的音乐创意，并在和弦选择中包含开弦。如果你想改变这种感觉，可以使用变调夹或者选择另一种调弦。}

\paragraph{Chords for FACGCE Tuning}
\eng{The approach to the following chords is twofold. First, I want to provide chords that capitalize on the sound of the tuning itself, making use of the open strings to create interesting open-string chord sounds within a key signature. Therefore, you will notice there is a chord for each note of the major scale. These chords sound wonderful clean or with light gain. To change key, apply a capo in the desired place. If using higher gain, you can also play the fretted notes and leave out the open strings for a tighter sound.}

\chn{以下和弦的做法有两个方面。首先，我想提供利用调弦本身声音的和弦，利用开弦来在调号内创建有趣的开弦和弦音色。因此，你会发现每个大调音阶音都有一个和弦。这些和弦在干音或轻失真时都非常好听。要改变调，只需将变调夹夹在所需位置。如果使用更强的失真，你也可以演奏按住的音符而省略开弦，以获得更紧致的声音。}

\eng{I also want to provide shapes that work well with higher gain so you do not need to learn two sets of chords. For example, for the Fmaj9 chord, play the top three strings if the full open shape is too muddy. With higher gain, these chords can become overly mushy if all open strings ring together, so be mindful of which strings are active.}

\chn{我还想提供适合更高失真的和弦形状，这样你就不需要学习两套和弦。例如，对于 Fmaj9 和弦，如果完整的开弦形状太混浊，可以只弹最顶部的三根弦。使用更高的失真时，如果所有开弦一起响，和弦可能会变得过于糊状，所以要注意哪些弦处于激活状态。}

\paragraph{FACGCE Chord Progressions in Context}
\eng{Below you will find chords and chord progressions from popular songs using FACGCE tuning. It is important to note that not all these chords are meant for strumming; some were originally played with picking patterns. However, they can still inspire your own ideas.}

\chn{下面你会看到使用 FACGCE 调弦的流行歌曲的和弦与和弦进行。需要注意的是，并非所有这些和弦都适合扫弦；有些最初是用拨奏模式演奏的。但它们仍然可以激发你的创作灵感。}

\eng{Examples include Covet – “Parachute”, Tiny Moving Parts – “You Lost Me”, Standards – “Smile”, American Football – “Never Meant”. Each progression demonstrates how the tuning can support open-string harmony and rich chord voicings.}

\chn{示例包括 Covet 的《Parachute》、Tiny Moving Parts 的《You Lost Me》、Standards 的《Smile》、American Football 的《Never Meant》。每个进行都展示了这种调弦如何支持开弦和声与丰富的和弦配位。}

\paragraph{FACGCE Arpeggios}
\eng{Below are major triad and minor triad arpeggio shapes for FACGCE tuning. The major shapes are based on F major, while the minor shapes are based on C minor. These shapes are movable, so you can adjust them based on the root of the major or minor chord you are using. Solid black notes on the charts show the root notes of each chord. Each chord has three arpeggios spanning an octave to cover all inversions: root position, first inversion, and second inversion.}

\chn{下面是 FACGCE 调弦的大小三和弦琶音形状。大调形状基于 F 大调，而小调形状基于 C 小调。这些形状是可移动的，因此你可以根据所用大调或小调和弦的根音进行调整。图表中的实心黑点表示每个和弦的根音。每个和弦都有三个跨越一个八度的琶音，以涵盖所有转位：根音形、第一转位和第二转位。}

\eng{As you study the charts, notice that the arpeggio shapes repeat at the octave. I have also marked where each arpeggio starts on each string. Although the root is positioned differently on each string, there are always three arpeggios: the root triad and the two possible inversions.}

\chn{学习图表时，注意琶音形状在八度时会重复。我还标出了每个琶音在各根弦上的起始位置。虽然根音在每根弦上的位置不同，但始终有三个琶音：根音三和弦和两种可能的转位。}

\paragraph{FACGCE Ideas in Context}
\eng{The following are five ideas I have written in FACGCE. I have included lots of noodling to get your brain working and to start noticing patterns and shapes available to you for writing licks in between your chords. Take your time with each one, and I recommend taking it phrase-by-phrase and building the idea until you can play it slowly. Examples two through five use a capo at the fourth fret.}

\chn{下面是我用 FACGCE 写的五个创意。我加入了大量即兴片段，让你的思维开始运转，并开始注意可用于在和弦之间写乐句的模式与形状。每个例子都慢慢练习，我建议逐句学习，直到你能慢速演奏。示例二到五使用第四品的变调夹。}

\paragraph{Chords for DAEAC\#E Tuning}
\eng{Just like FACGCE tuning, I want to share chords that take full advantage of the unique sound DAEAC\#E tuning offers. However, this tuning differs significantly from standard tuning, meaning it is much more difficult to find chords that align with every note in a key signature. So I take a slightly different approach by providing approachable chords and chord progressions that make the most of the open strings.} 

\chn{与 FACGCE 调弦一样，我也想分享充分发挥 DAEAC\#E 调弦独特声音的和弦。然而，这种调弦与标准调弦差异很大，意味着更难找到与调号中每个音符完全对齐的和弦。因此，我采用稍微不同的方法，提供易于上手的和弦和和弦进行，充分利用开弦。}

\eng{I also offer several chord progressions to illustrate how they can work together harmoniously. Try playing through these shapes and progressions to discover what you like. Then create your own progressions and explore riffs using the chord notes. These chords sound great cleanly or with a touch of light gain. If you want to use more distortion, you can play the fretted notes and exclude the open strings, but note that this may not work for every shape.}

\chn{我还提供了几个和弦进行，说明它们如何和谐地协同工作。尝试演奏这些形状和进行，以发现你喜欢的声音。然后创建你自己的进行，并探索使用和弦音的 riff。这些和弦在干音或轻失真下都很不错。如果你想使用更多失真，可以只弹按住的音符并省略开弦，但注意这并不适用于每个形状。}

\paragraph{DAEAC\#E Chords and Progressions}
\eng{The tuning is used for several ideas, including “For the Sake of Brevity Style Chord Progression Idea”, “Night of The Pencils Chord Progression Idea”, “Upbeat Midwest Emo Style Chord Progression”, and “Descending Style Chord Progression with A Twist”. Some of these use a capo on the fifth fret.} 

\chn{这种调弦用于多个创意，包括“取悦简洁风格和弦进行”、“铅笔之夜和弦进行创意”、“欢快中西部情绪摇滚风格和弦进行”，以及“带转折的下行风格和弦进行”。其中一些在第五品使用变调夹。}

\paragraph{DAEAC\#E Chord Progressions in Context}
\eng{The chords and progressions below are derived from popular songs using this tuning. Consider them a foundation and use the examples to spark more ideas. Examples include Tiny Moving Parts – “Always Focused”, “Caution”, and “Dakota”, as well as Algernon Cadwallader – “Spit Fountain”, Sleepy Dog – “Screamers”, and Mountains – “Stride”.}

\chn{下面的和弦与进行来源于使用这种调弦的流行歌曲。把它们视为基础，并用这些示例激发更多创意。示例包括 Tiny Moving Parts 的《Always Focused》、《Caution》和《Dakota》，以及 Algernon Cadwallader 的《Spit Fountain》、Sleepy Dog 的《Screamers》和 Mountains 的《Stride》。}

\paragraph{DAEAC\#E Arpeggios}
\eng{Here are major triad and minor triad arpeggios for DAEAC\#E tuning. The major shapes are in A major, and the minor shapes are in F\# minor. These shapes can be shifted based on the root of the chord you are using. Solid black notes indicate the root notes on each chart. Each chord has three arpeggios spanning an octave to cover root position and the two inversions.}

\chn{这里是 DAEAC\#E 调弦的大小三和弦琶音。大调形状是 A 大调，小调形状是 F\# 小调。你可以根据所用和弦的根音移动这些形状。实心黑点表示每个图表中的根音。每个和弦有三个跨越一个八度的琶音，覆盖根音形和两种转位。}

\paragraph{DAEAC\#E Ideas in Context}
\eng{The final ideas are written in DAEAC\#E and demonstrate how you can use the tuning to generate interesting melodic and chordal material. The examples are designed to get you exploring the unique voice of this tuning and to help you form ideas around the available open-string shapes.}

\chn{最后的创意是用 DAEAC\#E 写成的，展示了如何利用这种调弦生成有趣的旋律与和弦材料。这些示例旨在让你探索这种调弦的独特音色，并帮助你围绕可用的开弦形状形成创意。}
